\section{Conclusion}
The paper makes a strong case for the effectiveness of government spending in a liquidity trap, even if it warns against the fact that this effectiveness, measured by the marginal multiplier, may decline at a fast rate. It is important, however, not to forget the assumptions the model rests on. The taste shock that pushes the economy down to the zero lower bound, for example, implies a decrease of $23\%$ of GDP. Such a decline in output does not occur very often in real economies. For that reason, a very thorough assessment of the real economy should be done to check if the model can adequately inform policy-making for it.
\par
\bigskip
In general lines, we have been able to replicate the most important results of \citet{Erceg.2014}. Even when we could not get the exact values, our simulations behaved qualitatively in the same way as in the original paper. The results from the paper and from our replication are most different in the scenario with no inflation. For example, we got a multiplier that is roughly 0.1 smaller. In scenarios with Calvo parameter different than one, it seems we managed to get the exact same results of the paper (judging from the plots).
\par
\bigskip
The results of the model are very interesting and point to new research questions that could well be answered by refining the baseline model. Due to the importance of the Calvo parameter and the price setting framework in the model, alternative analysis with prices adjustments according to sticky information, rational inattention or state-dependent price settings are interesting for further research. Moreover, the simple New Keynesian model can be expanded to include a negative debt position in the steady-state, an open economy framework, the distinction between tax and debt financed government spending and an empirical comparison between the US and the euro area.
